\centerline{\bf \huge 9 класс}
\centerline{\normalsize{\it Работа рассчитана на \bf{180} минут}}

\problem{9.1}
Дан квадратный трёхчлен $f(x)$. Известно, что линейная функция $y= f(x+1)-f(x)$ обращается в ноль только при $x=2024$. При каком значении аргумента обращается в ноль функция $y=f(x-1)-f(x)$ ?

\problem{9.2}
На доске написаны четыре числа: $2,3,4$ и 9. За один шаг можно выбрать любые три из них, первое умножить на 2 , второе --- на 4 , а третье --- на 6 (при этом три старых числа стирают, а на их место записывают три новых). Можно ли через несколько шагов получить на доске четыре равных числа?

\problem{9.3}
Найдите значение суммы

$$
\begin{gathered}
S=\left(1^2+1 \cdot 3+3^2\right)+\left(3^2+3 \cdot 5+5^2\right)+\left(5^2+5 \cdot 7+7^2\right)+\cdots \\
+\left(97^2+97 \cdot 99+99^2\right)+\left(99^2+99 \cdot 101+101^2\right)
\end{gathered}
$$


\problem{9.4}
В треугольнике $A B C$ медиана $B M$ образует со стороной $B C$ прямой угол. На стороне $A B$ отмечена точка $X$ так, что $\angle B M C=\angle B M X$. Найдите, чему равно отношение $A X: X B$.

\problem{9.5}
На столе лежат $90$ карточек с числами от $1$ до $90$. Герман и Вова играют в следующую игру. Ходят по очереди. За один ход можно взять со стола любую карточку. Игра заканчивается, когда на столе останется две карточки. Вова выигрывает, если числа на оставшихся карточках отличаются ровно на $10$ . Иначе выигрывает Герман. Кто выигрывает при правильной игре, если начинает Герман?